% SCOPE: MERGED narrative arc. P1 RLVR optimizes one bit; nothing teaches
% handling own uncertainty (scale does not help). P2 direct repair backfires
% (self-distill suppresses hesitation; earlier calibration lowered
% overconfidence but never changed solution path) -> lesson sentence. P3 insight:
% pipelines reward FORM not FUNCTION (Gandhi as one clause). P4 reward causal
% effect (plain first; one sentence naming gold-vs-decoy shift; NO equations /
% teacher-forcing). P5 one-sentence positioning vs reference-likelihood line /
% AntiSD, three axes + twin clause. P6 preliminary but consistent (6/6 +
% hardest quartile) + one caveat. P7 four one-sentence contributions.
\section{Introduction}
\label{sec:intro}

Reinforcement learning with verifiable rewards, abbreviated RLVR, is now the
dominant recipe for improving mathematical reasoning in language models. The
method optimizes a single bit of feedback, namely whether the final answer is
correct \citep{shao2024deepseekmath,deepseekr1,zeng2025simplerlzoo,qwen3}. This
objective sharpens the model toward answers it can already reach. It teaches the
model nothing about handling its own uncertainty. Trained models stay
overconfident, commit early to a flawed approach, and rarely notice when their
own evidence has stopped supporting that approach
\citep{huang2024llmcannot,kadavath2022language}. The field's usual response is
scale: larger models, more compute, and longer traces. Scale raises the level of
problems a model can solve, but it does not change how the model behaves at the
edge of that level. Wherever the model's reliable reach ends, the same silent
overconfidence returns.

The direct repair for a broken uncertainty signal is to fix that signal, and
here the evidence is discouraging. Self-distillation trains a model on its own
verified solutions. Conditioning on a solution that is already known to be
correct suppresses the hesitation-bearing tokens, words like ``wait'' and ``let
me check'', that carry the model's deliberation. The traces then collapse into
short, confident paths, and reasoning degrades
\citep{kim2026selfdistill,opd2026rethinking,sdzero2026}. A different repair is to
make the model report its uncertainty honestly. In an earlier study on hard
mathematics we trained confidence as an explicit signal on top of a
metacognitive token layer. The intervention lowered overconfidence, but it did
not restore accuracy. The traces grew longer and more hesitant, yet the added
hesitation never turned into an actual change of solution path. The model
learned to say it was unsure without acting on that uncertainty. We carry
forward one lesson: calibration is necessary but not sufficient for
metacognitive control.

These outcomes share one cause: current pipelines reward the form of
self-reflection wherever it appears, but never its function. A model can learn to
produce verification-shaped text, such as ``let me double-check this step'', and
can even attach a well-calibrated confidence to it, while that text changes
nothing about the answer. No term in the objective asks whether the reflection
actually moved the model's evidence. \citet{gandhi2025cognitive} make a
compatible point from another angle: reinforcement learning amplifies cognitive
behaviors such as verification and backtracking only when those behaviors
already exist in the base model, so it selects among behaviors already present
rather than installing new ones. What is missing is not a better uncertainty
report. What is missing is training pressure that rewards a metacognitive act by
its causal effect on the model's evidence about the answer.

We build exactly that pressure. The model writes explicit metacognitive blocks
inside its chain of thought, opened by \metaopen{} and closed by \metaclose{},
and we reward each block by what it accomplishes. Concretely, we read the
model's own evidence about the answer just before a block and again just after
it, and we reward the block by how far it shifts that evidence away from a
plausible but wrong answer and toward the correct one. We read this shift under a
frozen reference model, contrasting a rule-based wrong answer, which we call a
decoy, against the gold answer, and we name the resulting reward \pmishift{}.
Correctness stays the primary objective, and \pmishift{} is an auxiliary shaping
term.

This reward sits in a recent line of work that scores a reasoning trace by the
likelihood it assigns to a reference answer
\citep{yu2025rlpr,zhang2025crm,mulgund2026rlsd}. Its closest neighbor is
AntiSD~\citep{shen2026antisd}, from which we differ by rewarding a
gold-versus-decoy shift localized to an explicit metacognition span and validated
with a causal twin.

Our results are preliminary but consistent. On held-out mathematics the
metacognitive model beats its matched twin on all six benchmark comparisons, and
its advantage concentrates on the hardest problems, reaching $+34.8$ percentage
points in the top difficulty quartile even though the model emits metacognitive
blocks at nearly the same rate across every difficulty. These numbers come from a
single seed and from a reward package that bundles \pmishift{} with several
auxiliary terms, so they measure the package rather than \pmishift{} in
isolation, and the decomposition that would isolate it is ongoing.

\paragraph{Contributions.}
\begin{enumerate}
\item An \textbf{insight and its operationalization}: calibration is necessary
  but not sufficient for metacognitive control, so the signal worth rewarding is
  the causal effect of a metacognitive act on the model's evidence, which we turn
  into an executable definition of in-reasoning metacognition as explicit
  \metaopen{}\dots\metaclose{} blocks that govern when to verify and when to
  redirect.
\item \textbf{The \pmishift{} reward}: we score each metacognitive block by how
  far it shifts a gold-versus-decoy evidence margin across the span it occupies,
  a signal that differs from a plain per-token answer-likelihood reward by being
  contrastive, span-localized, and tied to the metacognitive region.
\item A \textbf{matched-base causal design and analysis suite}: we train a twin
  model that is identical except for the removed metacognitive content and
  byte-audited to share every other setting, and we add decomposition, placebo,
  difficulty-stratification, and calibration analyses that attribute any gain to
  the mechanism.
\item A \textbf{preliminary empirical characterization} on math reasoning with
  Qwen3-8B, showing that the metacognitive model wins all six held-out
  comparisons and that its advantage hides in pooled averages because it lives
  almost entirely in the hardest problems, so any comparison that mixes easy and
  hard problems understates the effect and must be split by difficulty.
\end{enumerate}
